\section{Tool Procurement}

Before the analysis begins, each role must contribute to a documented tool selection decision. Three capability areas have been identified as immediate requirements for this case.

\subsection{Required Capability Areas}

\textbf{Area 1: Cloud API Interaction}\\
Tools must allow listing, retrieving, and hashing objects from a cloud object store via command
line, without requiring direct infrastructure access.

\textbf{Area 2: Structured Log Analysis}\\
Tools must parse and query JSON-formatted event logs, filtering by event type, timestamp, and
actor.

\textbf{Area 3: Local Cloud Emulation}\\
A capability to reproduce a minimal cloud object-storage API locally, for training purposes,
without incurring cost or requiring a real provider account.

\subsection{Background: Object Storage and S3}

Before listing the tools, it is worth clarifying two things that later sections assume you
already understand.
\medskip

\noindent
\textbf{What ``S3'' actually is.} S3 (Simple Storage Service) is the name of Amazon's object
storage product, but over time its HTTP API has become a de-facto \emph{industry protocol}:
``S3-compatible'' storage is now offered by most cloud providers and by open-source software,
regardless of who actually built it. Unlike a filesystem, an object store has no directories in
the traditional sense: every object is stored flat inside a \emph{bucket}, addressed by a
\emph{key} (its name/path string) and accompanied by \emph{metadata} (size, content type,
last-modified time, and optionally custom key-value tags). Whenever this laboratory says
``interact with the bucket'', it means issuing HTTP requests against this API, normally
through the \texttt{aws cli}, which simply builds and sends those requests for you.
\medskip

\noindent
\textbf{Why we interact with S3 directly at all.} NovaDynamics' actual cloud provider is, of
course, not accessible to us: we only have the exports listed in Section 1.4. Section 4 of
this laboratory has you interact with a \emph{local stand-in} bucket instead. This is not
redundant busywork: it is the only way, short of a real provider account, to see first-hand what
the S3 API actually returns for an operation like ``list objects'' or ``get object metadata''.
That first-hand experience is what lets you later judge, in Sections 5--6, whether a given export
is complete, plausible, or has gaps.

\subsection{Background: Docker}

Docker lets you run a piece of software (here, MinIO) inside an isolated
\emph{container}, without installing it directly onto your VM's operating system and without it
interfering with anything else on the system. You do not need to understand Docker's internals
for this laboratory; you only need three operations, all of which are shown step by step in
Section 3: \emph{run} a container (start the software), \emph{list} running containers (check
it is alive), and \emph{stop/remove} a container (clean up afterwards). We use Docker here purely
as a convenient, disposable way to stand up a local S3-compatible server, \emph{not} as a
forensic subject in its own right (container forensics is a distinct topic, not covered by this
laboratory).

\subsection{Tool Versions and Reproducibility}

Every command in Section 3 specifies an explicit version for the tools it installs or runs,
rather than pulling whatever the current release happens to be. This is deliberate, and it is a
forensic requirement rather than a stylistic preference: an analysis that cannot be reproduced is
an analysis that cannot be defended. Different versions of the same tool can behave differently
on identical input, and the versions recorded here are the ones under which this laboratory was
verified end to end.
\medskip

\noindent
This is not a hypothetical concern. Two concrete examples encountered while preparing this
laboratory:

\begin{itemize}[leftmargin=1.2cm]
	\item An earlier draft used LocalStack as the local emulator. Since March 2026, its Docker
	      image requires a free account and an auth token even to start, which defeats the
	      ``no account, no cost'' constraint of Area 3 above. Pinning to an older release did not
	      solve the problem either, because recent versions of \texttt{aws cli} (from v2.23,
	      January 2025) send an additional checksum header on every upload that older emulator
	      releases reject.
	\item MinIO itself is therefore pinned to a specific release tag below. If a future release
	      changes its licensing or its API behaviour, the pinned tag keeps this laboratory
	      reproducible until it is formally revised.
\end{itemize}

\noindent
Record in your notes the exact version of every tool you actually used, including \texttt{aws cli}
(\texttt{aws -\/-version}) and \texttt{jq} (\texttt{jq -\/-version}), even where this text does not
pin them: your report should let a third party re-run your analysis and obtain the same results.

\subsection{Pre-approved Tools}

\begin{description}[leftmargin=1.5cm, style=nextline]
	\item[\texttt{aws cli}] Standard command-line interface for interacting with S3-compatible
	      APIs; usable against both real endpoints and local emulators.
	\item[\texttt{jq}] Command-line JSON processor, for querying and filtering the CloudTrail-style
	      log export.
	\item[MinIO (Docker)] Local, open-source, cost-free emulation of an S3-compatible object
	      store; run inside the same CAINE/Kali VM already used in previous laboratories.\\
		  Pinned to release \texttt{RELEASE.2025-09-07T16-13-09Z}, the version under which this
		  laboratory was verified.
	\item[\texttt{mc} (MinIO Client)] Companion command-line client for MinIO, used in Section 4.5
	      to configure event notifications, a capability not exposed through the plain S3 API.
	\item[\texttt{sha256sum}] Already used in Laboratories 2 and 3, reused here for verifying
	      object integrity after acquisition.
\end{description}
